$xy$平面上の任意の点を$X(x,y)$とする。
線分$PQ$上の点は、$0\leqq t\leqq1$を用いて端点を補間して表せるから、
\[
\begin{aligned}
X(x,y)\in D
&\Longleftrightarrow
\begin{gathered}
\text{「}\;
\left\{
\begin{aligned}
0&\leqq\theta\leqq\dfrac{\pi}{2}\\
0&\leqq t\leqq1\\
x&=8t\cos\theta\\
y&=(1-t)\sin\theta
\end{aligned}
\right.\\
\text{を満たす実数 }\theta,t\text{ が存在する」}
\end{gathered}
\end{aligned}
\]

$0<x<8$を固定し、
\[
\cos\alpha=\frac{x}{8},
\qquad
0<\alpha<\frac{\pi}{2}
\]
とする。$x=8t\cos\theta>0$であるから、$t>0$かつ
$\cos\theta>0$であり、
\[
t=\frac{x}{8\cos\theta}
\]
となる。$\cos\theta$は
$0\leqq\theta\leqq\dfrac{\pi}{2}$で減少するので、
\[
\begin{aligned}
0\leqq t\leqq1
&\Longleftrightarrow
0<\frac{x}{8\cos\theta}\leqq1\\
&\Longleftrightarrow
x\leqq8\cos\theta\\
&\Longleftrightarrow
\cos\alpha\leqq\cos\theta\\
&\Longleftrightarrow
0\leqq\theta\leqq\alpha
\end{aligned}
\]
である。また、
\[
f_x(\theta)
=\sin\theta-\frac{x}{8}\tan\theta
\]
とおく。ここで、所属条件から$t$を消去する。
\[
\begin{aligned}
X(x,y)\in D
&\Longleftrightarrow
\begin{gathered}
\text{「}\;
\left\{
\begin{aligned}
0&\leqq\theta\leqq\dfrac{\pi}{2}\\
0&\leqq t\leqq1\\
x&=8t\cos\theta\\
y&=(1-t)\sin\theta
\end{aligned}
\right.\\
\text{を満たす実数 }\theta,t\text{ が存在する」}
\end{gathered}\\
&\Longleftrightarrow
\begin{gathered}
\text{「}\;
\left\{
\begin{aligned}
0&\leqq\theta\leqq\alpha\\
t&=\dfrac{x}{8\cos\theta}\\
y&=\left(1-\dfrac{x}{8\cos\theta}\right)\sin\theta
\end{aligned}
\right.\\
\text{を満たす実数 }\theta,t\text{ が存在する」}
\end{gathered}\\
&\Longleftrightarrow
\begin{gathered}
\text{「}\;
\left\{
\begin{aligned}
0&\leqq\theta\leqq\alpha\\
y&=f_x(\theta)
\end{aligned}
\right.\\
\text{を満たす実数 }\theta\text{ が存在する」}
\end{gathered}
\end{aligned}
\]

残った$\theta$を消去するため、$f_x$の値域を求める。

$0\leqq\theta\leqq\alpha<\dfrac{\pi}{2}$では
\[
f_x'(\theta)
=\cos\theta-\frac{x}{8\cos^2\theta}
=\frac{\cos^3\theta-\dfrac{x}{8}}{\cos^2\theta}
\]
である。ここで、
\[
\cos^3\beta=\frac{x}{8},
\qquad
0<\beta<\frac{\pi}{2}
\]
となる$\beta$をとる。$0<\dfrac{x}{8}<1$であるから、
\[
\frac{x}{8}<\left(\frac{x}{8}\right)^{1/3}<1
\]
である。$\cos\alpha=\dfrac{x}{8}$、
$\cos\beta=\left(\dfrac{x}{8}\right)^{1/3}$より
\[
\cos\alpha<\cos\beta<1
\]
であり、$0\leqq\theta\leqq\dfrac{\pi}{2}$では$\cos\theta$が減少するので、
\[
0<\beta<\alpha
\]
となる。

この範囲では$\cos^2\theta>0$である。また、
$\cos\theta$は減少するから、$f_x'(\theta)$の符号は
\[
\left\{
\begin{aligned}
f_x'(\theta)&>0 &&(0\leqq\theta<\beta)\\
f_x'(\beta)&=0\\
f_x'(\theta)&<0 &&(\beta<\theta\leqq\alpha)
\end{aligned}
\right.
\]
となる。このとき、
\[
\begin{aligned}
f_x(\beta)
&=\sin\beta-\cos^3\beta\tan\beta\\
&=(1-\cos^2\beta)\sin\beta\\
&=\left\{1-\left(\frac{x}{8}\right)^{2/3}\right\}^{3/2}
\end{aligned}
\]
であり、$f_x(0)=f_x(\alpha)=0$である。したがって、増減表は
\[
\begin{array}{c|ccccc}
\theta&0&&\beta&&\alpha\\ \hline
f_x'(\theta)&&+&0&-&\\ \hline
f_x(\theta)&0&\nearrow&
\left\{1-\left(\dfrac{x}{8}\right)^{2/3}\right\}^{3/2}
&\searrow&0
\end{array}
\]
となる。この表より、$f_x(\theta)$は$\theta=\beta$のとき最大となり、最大値は
\[
\left\{1-\left(\frac{x}{8}\right)^{2/3}\right\}^{3/2}
\]
である。さらに、$f_x$は$0\leqq\theta\leqq\alpha$で連続であるから、$f_x$の値域は
\[
0\leqq f_x(\theta)\leqq
\left\{1-\left(\frac{x}{8}\right)^{2/3}\right\}^{3/2}
\]
である。よって、$0<x<8$を固定すると、
\[
\begin{aligned}
X(x,y)\in D
&\Longleftrightarrow
\begin{gathered}
\text{「}\;
\left\{
\begin{aligned}
0&\leqq\theta\leqq\alpha\\
y&=f_x(\theta)
\end{aligned}
\right.\\
\text{を満たす実数 }\theta\text{ が存在する」}
\end{gathered}\\
&\Longleftrightarrow
0\leqq y\leqq
\left\{1-\left(\frac{x}{8}\right)^{2/3}\right\}^{3/2}
\end{aligned}
\]
となる。

初めの補間条件から、$x=0$では$0\leqq y\leqq1$であり、
$x=8$では$t=1$かつ$\theta=0$となるので$y=0$である。これらを統合すると、
\[
\begin{aligned}
X(x,y)\in D
&\Longleftrightarrow
\left\{
\begin{aligned}
0&\leqq x\leqq8\\
0&\leqq y\leqq
\left\{1-\left(\dfrac{x}{8}\right)^{2/3}\right\}^{3/2}
\end{aligned}
\right.
\end{aligned}
\]

この領域$D$を$x$軸のまわりに回転した立体の体積$V$は
\[
V
=\pi\int_0^8
\left\{1-\left(\frac{x}{8}\right)^{2/3}\right\}^{3}
\,dx
\]
である。$x=8u^3$とおくと、$0\leqq u\leqq1$かつ$dx=24u^2\,du$であるから、
\[
\begin{aligned}
V
&=24\pi\int_0^1u^2(1-u^2)^3\,du\\
&=24\pi\int_0^1
\left(u^2-3u^4+3u^6-u^8\right)\,du\\
&=24\pi
\left(\frac13-\frac35+\frac37-\frac19\right)\\
&=\frac{128\pi}{105}
\end{aligned}
\]
